% SBSI -- preliminary status slides.
% Build:  pdflatex sbsi_slides.tex   (run from this directory)
\documentclass[aspectratio=169,10pt]{beamer}

\usetheme{default}
\usecolortheme{dove}
\usepackage[T1]{fontenc}
\usepackage{helvet}
\usepackage{amsmath,amssymb,bm}
\usepackage{graphicx}
\usepackage{booktabs}
\usepackage{tabularx}
\usepackage{array}
\newcolumntype{L}[1]{>{\raggedright\arraybackslash}p{#1}}
\newcolumntype{Y}{>{\raggedright\arraybackslash}X}
\renewcommand{\familydefault}{\sfdefault}

\setbeamertemplate{navigation symbols}{}
\setbeamertemplate{footline}{%
  \hfill\usebeamerfont{footline}\color{gray!80}\scriptsize\insertframenumber\hspace{4mm}\vspace{2mm}}
\setbeamerfont{frametitle}{size=\large,series=\bfseries}
\setbeamercolor{frametitle}{fg=black}
\setbeamertemplate{frametitle}{\vspace{3mm}\insertframetitle\par
  \vspace{-1mm}{\color{gray!50}\rule{\textwidth}{0.4pt}}}
\setbeamertemplate{itemize item}{\color{gray!70}\textbf{--}}
\setbeamertemplate{itemize subitem}{\color{gray!70}\textbf{\textperiodcentered}}
\setbeamercolor{block title}{fg=black,bg=black!6}
\setbeamercolor{block body}{bg=black!3}

\graphicspath{{../figures/}}

\newcommand{\xh}{\hat{\mathbf{x}}}
\newcommand{\xt}{\mathbf{x}}
\newcommand{\nb}{\mathbf{n}}
\newcommand{\Rsf}{R_{\rm self}}
\newcommand{\Rbl}{R_{\rm blend}}

\begin{document}

% ----------------------------------------------------------------- 1
\begin{frame}{SBSI --- simulation-based shear inference for blended, selected galaxies}
\vspace{2mm}
{\large Preliminary status, \today}
\vspace{4mm}

\begin{itemize}\itemsep5pt
  \item \textbf{Problem.} Weak-lensing shear from an LSST-like catalogue where galaxies
        \emph{overlap} (blending) and where \emph{cuts} are applied. Target: multiplicative
        bias $|m|<0.3\%$.
  \item \textbf{Approach.} Not a calibration. Learn the measurement likelihood
        $p(\xh\mid\xt,\nb)$ with a conditional normalizing flow trained on image
        simulations (ngmix shapes), then do inference through it.
  \item \textbf{Shear enters the prior}, not the data: it is a transformation $S_\gamma$ of the
        \emph{scene}. The estimator is a score / Fisher ratio (BFD-like). No true galaxy
        property is ever used at inference time.
  \item \textbf{Two response channels.} $R=\Rsf+\Rbl$: the galaxy's own shape response, and
        the coupling to its neighbours. $\Rbl$ currently comes from an external
        (BlendEMU) emulator; folding it into the flow is the next step.
\end{itemize}
\end{frame}

% ----------------------------------------------------------------- 2
\begin{frame}{The model, and its two derivatives \quad{\normalsize\color{gray}(MATH.md \S6, steps 1--4)}}
\vspace{1mm}
{\small A scene $\theta=(\xt,\nb)$ --- primary plus neighbours --- is drawn from a prior $p_0$, and the
flow turns it into a measurement. Shear acts on the \textbf{prior}, pushing it forward along
$S_\gamma$. Keeping the detection classifier, the measured density of one object is}
\vspace{-2mm}
\[ A(\xh\mid\gamma)\;=\;\int d\theta\;
   \underbrace{p_{\rm flow}(\xh\mid\theta)}_{\text{flow}}\;
   \underbrace{P_{\rm det}(\theta)}_{\text{classifier}}\;p_\gamma(\theta),
   \qquad\quad p_\gamma\;=\;S_\gamma\#p_0 . \]

\vspace{-2mm}
{\small \textbf{The generator.} $p_\gamma$ is the only factor carrying $\gamma$. A pushforward obeys a
continuity equation, which gives its logarithmic derivative in closed form:}
\vspace{-2mm}
\[ u\;\equiv\;\partial_\gamma\log p_\gamma\big|_0
   \;=\;-\big(v\cdot\nabla\log p_0+\nabla\!\cdot\!v\big),
   \qquad\quad v\;\equiv\;\partial_\gamma S_\gamma\big|_0 . \]

\vspace{-2mm}
{\small \textbf{Two derivatives of $\log A$} --- Fisher's identity, then Louis's --- both as averages
over a node bank $\theta_k\sim p_0$ with posterior weights
$w_k\propto p_{\rm flow}(\xh_i\mid\theta_k)P_{\rm det}(\theta_k)$:}
\vspace{-2mm}
\[ \underbrace{s_i\;=\;\partial_\gamma\log A\;=\;\mathbb{E}_{w_i}[u]}_{\text{score}},
   \qquad\qquad
   \underbrace{\mathcal{I}_i\;=\;-\partial^2_\gamma\log A
   \;=\;-\mathbb{E}_{w_i}\big[\partial_\gamma u\big]-\mathrm{Var}_{w_i}(u)}_{\text{information}} . \]

\vspace{-1mm}
{\scriptsize\color{black!75}
The flow is \emph{never} differentiated with respect to $\gamma$ --- $\gamma$ is not one of its
inputs. And $u$ does not depend on the galaxy: $u_k$ and $\partial_\gamma u_k$ are built \emph{once
per node} and shared across the whole catalogue, leaving one flow call per (object, node). The price
is $\nabla\log p_0$ over the whole scene.}
\end{frame}

% ----------------------------------------------------------------- 3
\begin{frame}{Selection, the estimator, and $\Rbl$ \quad{\normalsize\color{gray}(MATH.md \S6, steps 5--7)}}
\vspace{1mm}
{\small \textbf{Selection.} A kept object is drawn from $A$ restricted to the cut and
\emph{renormalised} --- one normaliser handles detection and the cut together:}
\vspace{-2mm}
\[ \log p_{\rm keep}\;=\;\log W(\xh)+\log A(\xh\mid\gamma)-\log P(\gamma),
   \qquad P(\gamma)=\mathbb{E}_{p_\gamma}\big[\Pi\big],\quad \Pi=P_{\rm pass}P_{\rm det} . \]

\vspace{-2mm}
{\small $W$ is $\gamma$-free and $=1$ for a galaxy in hand, so differentiating once and twice subtracts
the \emph{same two population numbers} from \emph{both} moments --- same node bank, reweighted by prior
survival $\Pi_k$:}
\vspace{-2mm}
\[ \langle s\rangle_{\rm sel}=\mathbb{E}_{\Pi}[u],
   \qquad\qquad
   \mathcal{I}_{\rm sel}=-\mathbb{E}_{\Pi}\big[\partial_\gamma u\big]-\mathrm{Var}_{\Pi}(u) . \]

\vspace{-1mm}
{\small One Newton step on the kept log-likelihood --- slope over curvature --- is the estimator:}
\vspace{-1mm}
\begin{center}
\fbox{$\displaystyle\;\hat\gamma\;=\;-\frac{\ell'(0)}{\ell''(0)}
   \;=\;\frac{\sum_i s_i-N\langle s\rangle_{\rm sel}}
             {\sum_i\mathcal{I}_i-N\,\mathcal{I}_{\rm sel}}\;$}
\end{center}

\vspace{0mm}
{\footnotesize \textbf{Cut vs.\ detection.} A cut on \emph{measured} $\xh$ is a function of data already
held, so it cancels from the per-object posterior: $P_{\rm pass}$ appears \emph{only} in the population
term. Detection is \emph{not} determined by $\xh$, so $P_{\rm det}$ survives in both.}

\vspace{1mm}
{\footnotesize \textbf{Where $\Rbl$ enters.} With $\nb$ latent and $S_\gamma$ acting on the whole scene,
$u$ carries neighbour components --- separation, neighbour shape and size, all spin-2 --- so
$s_i=\sum_kw_ku_k$ picks the blend response up with \emph{no separate term}.}

\vspace{0mm}
{\scriptsize\color{black!75}
It vanishes when the flow is blind to the neighbour position angle --- that angle's posterior equals
its prior, and isotropy kills the sum. Today's conditioning is scalar, so $\Rbl$ is still supplied
externally: $m=R_{\rm sim}/(\Rsf+\Rbl)-1$.}
\end{frame}

% ----------------------------------------------------------------- 4
\begin{frame}{Trained models}
\vspace{2mm}
{\footnotesize\setlength{\tabcolsep}{4pt}\renewcommand{\arraystretch}{1.15}
\begin{tabularx}{\textwidth}{@{}L{2.3cm}L{2.6cm}YY@{}}
\toprule
\textbf{Model} & \textbf{Provides} & \textbf{Architecture} & \textbf{Training} \\
\midrule
Measurement flow\newline\emph{V2 \texttt{dom6x6}} &
$p_{\rm flow}(\xh|\xt,\nb)$: 4D output $(g_1,g_2,\mathrm{mag},\log R_{50})$ from 8 scene features &
conditional mean-affine flow, 10 blocks, hidden 256, 3 layers, SiLU, 16D context &
3.4M/0.6M rows, 80 epochs, batch 8192, lr $7\!\times\!10^{-4}$; NLL $+\ \lambda_r$ response penalty
$+\ \lambda_\theta$ spin-2 coupling; SWA epochs 73--80; \textbf{16 seeds} \\[2pt]
Detection classifier\newline\emph{response-aware} &
$P_{\rm det}(\xt,\nb)$ &
MLP $4\times256$, SiLU, 28 inputs incl.\ pair-frame $\cos2\varphi,\sin2\varphi$ &
BCE $+\ \lambda=300$ regularisation onto the measured detection response \\[2pt]
Blend emulator\newline\emph{BlendEMU, tuned} &
$\Rbl$ &
external lookup, joined per object &
valid only in its box: true $\mathrm{mag}<26$, $0.3<R_e<1.5''$ \\
\bottomrule
\end{tabularx}}

\vspace{1mm}
{\scriptsize\color{black!75}Constant-shear population, $N=11{,}674{,}408$ in-domain objects:\;
$R_{\rm sim}=0.8605$,\; $R_{\rm flow}=0.7258$,\; $\Rbl=0.1358$\; $\Rightarrow$\;
16-seed ensemble $m=-0.12\pm0.15\%$ (seed scatter $0.61\%$).}

\end{frame}

% ----------------------------------------------------------------- 5
\begin{frame}{Training, and the spread over seeds}
\vspace{2mm}
\begin{columns}[t]
\begin{column}{0.5\textwidth}\centering
\includegraphics[width=\linewidth]{fid_fig1_seed_loss.png}\\[2mm]
{\scriptsize\color{black!75}Validation NLL per epoch, all 16 seeds. Plateau at $1.429$; the
checkpoint used is an SWA average over epochs 73--80.}
\end{column}
\begin{column}{0.5\textwidth}\centering
\includegraphics[width=0.96\linewidth]{fid_fig3_bias_true_neighbours.png}\\[2mm]
{\scriptsize\color{black!75}Multiplicative bias per seed. Scatter between seeds is $0.61\%$, so the
16-seed ensemble gives $m=-0.12\pm0.15\%$ --- inside the $\pm0.3\%$ target.}
\end{column}
\end{columns}
\end{frame}

% ----------------------------------------------------------------- 6
\begin{frame}{Self response: the flow alone, against half-shear truth}
\vspace{-1mm}
\begin{center}
\includegraphics[width=0.94\textwidth]{fid_fig5_selfresp_halfshear.png}
\end{center}
\vspace{-2mm}
{\footnotesize\color{black!75}
Half-shear simulation legs isolate $\Rsf$, so the emulator has no freedom here --- this is the flow
on its own. It tracks the truth to within about $\pm2\%$ across primary flux, size and neighbour
flux, degrading at the bright end ($-7\%$ in the last S/N bin) and at the crowded tail.
Band $=$ spread over the 16 seeds; $N=2{,}364{,}527$.}
\end{frame}

% ----------------------------------------------------------------- 7
\begin{frame}{Full response: flow $+$ emulator, against the constant-shear simulation}
\vspace{-1mm}
\begin{center}
\includegraphics[width=0.94\textwidth]{fid_fig2_response_vs_properties.png}
\end{center}
\vspace{-2mm}
{\footnotesize\color{black!75}
$R_{\rm flow}+\Rbl$ vs.\ the simulated response. Residuals stay within roughly $\pm5\%$ per bin and
average to $m=-0.12\pm0.15\%$ over the population. The population is the emulator's in-domain box
(true $\mathrm{mag}<26$, $R_e>0.3''$); rows outside it are dropped, not zero-filled.}
\end{frame}

% ----------------------------------------------------------------- 8
\begin{frame}{Selection: response under near-domain \emph{measured} cuts}
\vspace{-1mm}
\begin{center}
\includegraphics[height=5.5cm]{fid_fig4_selection_near_domain.png}
\end{center}
\vspace{-2mm}
{\footnotesize\color{black!75}
Cuts move the mean response by up to $17\%$ (left); the model follows most of that move. Residual
bias at the cut (right): $-0.12\pm0.15\%$ no cut,\; $+0.20\pm0.16\%$ at $\mathrm{mag}<26$,\;
$+0.63\pm0.17\%$ at $\mathrm{mag}<25$,\; $+0.66\pm0.17\%$ at $\mathrm{S/N}>10$ --- and
$+4.46\pm0.14\%$ at $R>0.70''$, the one clear failure.
\textbf{Open:} $\Rbl$ still external; shear also moves galaxy \emph{positions} (not yet modelled);
all of the above is in-domain only.}
\end{frame}

\end{document}
